\section{The dyadic-sum source and its weighted operator}
\label{sec:source}

Let $G_2$ be the undirected looped graph with vertex set $\N$ and
\begin{equation}
  m\sim n \quad\Longleftrightarrow\quad
  m+n=2^a\text{ for some integer }a\ge1.
  \label{eq:edge-relation}
\end{equation}
Its one-sided edge shift consists of sequences
$(n_j)_{j\ge0}$ satisfying $n_j\sim n_{j+1}$.  One edge is one unit of
time, and the formal marker $z$ records that time.  A primitive orbit, when
needed, is a cyclic vertex word of least shift period, modulo rotation.
Dyadic edge labels are constraints derived from a vertex word; they are not
the primitive objects.

For $s\in\C$, define the coefficient array $h_s$ by
\cref{eq:intro-matrix}, with
$n^{-s/2}=\exp[-(s/2)\log n]$ and the real logarithm.  We say that $h_s$
has a \emph{bounded realization} if there is $H_s\in\Bop(\ell^2(\N))$ with
matrix coefficients $\ip{H_se_n}{e_m}=h_s(m,n)$.  This formulation matters
when $\sigma\le0$: a column then need not lie in $\ell^2$, so we do not
silently introduce an unbounded operator on finitely supported vectors.  On
odd vertices, $A_s$ denotes the corresponding compression whenever the
bounded realization exists.

\begin{table}[t]
  \centering
  \caption{The typed objects used in the operator and walk ledgers.}
  \label{tab:typed-objects}
  \small
  \begin{tabular}{@{}lll@{}}
    \toprule
    Object & Mathematical type & Not identified with \\
    \midrule
    $n$ & positive-integer vertex & rational-prime primitive \\
    $q_i=2^{a_i}$ & derived edge label & temporal orbit \\
    $(n_1,\ldots,n_r)$ & based closed vertex walk & unordered labels \\
    $z$ & one-edge time marker & valuation weight \\
    $2^{-krs}$ & $r$-step block weight & time marker \\
    $H_s$ & bounded weighted adjacency, when legal & nonreal self-adjoint operator \\
    \bottomrule
  \end{tabular}
\end{table}

\begin{lemma}[Left--right phase factorization]
\label{lem:phase}
Let $s=\sigma+it$, let $U_t e_n=n^{-it/2}e_n$, and write
$H_s^{(N)}=(h_s(m,n))_{1\le m,n\le N}$ and
$U_t^{(N)}=\diag(n^{-it/2})_{1\le n\le N}$.  Entrywise, and hence for every
finite matrix,
\begin{equation}
  H_s^{(N)}=U_t^{(N)}H_\sigma^{(N)}U_t^{(N)}.
  \label{eq:finite-phase}
\end{equation}
If the coefficient arrays have bounded realizations, then
$H_s=U_tH_\sigma U_t$.  Consequently boundedness, compactness, singular
values, Schatten membership, and the corresponding norms depend only on
$\sigma$.
\end{lemma}

\begin{proof}
The $(m,n)$ coefficient on the right of \cref{eq:finite-phase} is
$m^{-it/2}h_\sigma(m,n)n^{-it/2}=h_s(m,n)$.  Each finite compression is an
ordinary matrix identity.  If $H_\sigma$ is bounded, the same calculation
gives the full bounded identity.  Conversely, $U_t^*H_sU_t^*$ is a bounded
operator with coefficients $h_\sigma$.  Left and right multiplication by
unitaries preserves singular values and all properties listed above.
\end{proof}

The factorization is not a unitary conjugacy: the two copies of $U_t$ point
in the same direction.  It therefore transfers neither spectra nor powers,
traces, or determinants.  All such identities below retain the complex
parameter and come instead from an actual basis-reordering direct sum.

The loop locations and valuation constraint are elementary but structural.
There is a loop at $m$ precisely when $2m=2^a$, hence precisely at
$m=2^k$, $k\ge0$.

\begin{lemma}[Valuation preservation]
\label{lem:valuation}
If $m+n=2^a$ with $m,n\in\N$, then $\vTwo(m)=\vTwo(n)$.  If this common
value is $k$ and $m=2^ku$, $n=2^kv$, then $u,v$ are odd and
$u+v=2^{a-k}$.
\end{lemma}

\begin{proof}
If the two valuations differed, the valuation of their sum would equal the
smaller one.  Because $m,n<2^a$, that value is strictly smaller than $a$,
contradicting $\vTwo(m+n)=a$.  Dividing the equality by $2^k$ gives the
second statement.
\end{proof}
